%% Subdirect representation for nonempty finitary algebras.

\subsection{Subdirect products and irreducible quotients}
\label{sec:p10-subdirect-products}

Fix a finitary signature \(\tau\); its set of operation symbols need not
be finite.  All algebras have nonempty carriers, and all product indices
form sets.  We work in set theory with the axiom of choice; its use
through Zorn's lemma in the representation theorem will be explicit.
The source for this toolkit is Chapter~II, \S8 of the 2012 Millennium
Edition of \textcite{BurrisSankappanavar2012}, with the quotient
correspondence from Chapter~II, Theorem~6.20.

Write \(\operatorname{Con}A\) for the congruences on \(A\), ordered by
inclusion, and put
\[
 \Delta_A=\{(a,a):a\in A\},\qquad \nabla_A=A\times A.
\]
Intersections of congruences are congruences; the empty intersection is
\(\nabla_A\).  Thus a family has meet given by intersection and join
given by the intersection of all its common congruence upper bounds.
In particular, \(\theta\vee\phi\) is the least congruence containing
both \(\theta\) and \(\phi\).  The \emph{principal congruence}
\(\operatorname{Cg}_A(a,b)\) is the intersection of all congruences
containing \((a,b)\).  The quotient map is
\(q_\theta:A\to A/\theta\), \(a\mapsto[a]_\theta\), and
\(\ker h=\{(a,b):h(a)=h(b)\}\) for a homomorphism \(h\).

\begin{hbdefinition}{def:p10-subdirect-product}{Subdirect products and embeddings}
For an indexed family \((A_i)_{i\in I}\) of \(\tau\)-algebras, its
direct product has the Cartesian product as carrier and coordinatewise
operations:
\[
 f^{\prod_{i\in I}A_i}(x_1,\ldots,x_n)(i)
       =f^{A_i}(x_1(i),\ldots,x_n(i)).
\]
This includes nullary operations.  For \(I=\varnothing\), the product
is the one-element \(\tau\)-algebra whose element is the empty tuple.
For \(i\in I\), let \(\pi_i\) be the coordinate projection.

A subalgebra \(B\leq\prod_{i\in I}A_i\) is a nonempty subset closed
under the basic operations, with those operations restricted to \(B\).
It is a \emph{subdirect product} when every restricted projection
\(\pi_i|_B:B\to A_i\) is surjective.
A \emph{subdirect embedding} is an injective homomorphism
\(e:A\to\prod_{i\in I}A_i\) whose image is a subdirect product.
Equivalently, \(e\) is injective and every \(\pi_i\circ e\) is onto.
The image may be a proper subalgebra of the full product.
For an empty index set a subdirect product is exactly the one-element
product, and an algebra embeds subdirectly into it exactly when that
algebra is trivial, meaning that its carrier has one element.
\index{subdirect product}\index{subdirect embedding}
\end{hbdefinition}

\begin{hbtheorem}{thm:p10-subdirect-kernels}{proved=annex:p10.subdirect.1}{The congruence-intersection criterion}
Let \(h_i:A\to A_i\) be homomorphisms, for \(i\in I\), and define
\[
 e:A\longrightarrow\prod_{i\in I}A_i,\qquad
 e(a)=(h_i(a))_{i\in I}.
\]
Then \(e\) is a homomorphism and
\[
 \ker e=\bigcap_{i\in I}\ker h_i.
\]
It is an embedding exactly when that intersection is \(\Delta_A\);
if each \(h_i\) is onto, this is exactly the condition for \(e\) to
be subdirect.

In particular, for congruences \(\theta_i\in\operatorname{Con}A\),
put \(\theta=\bigcap_{i\in I}\theta_i\).  The map
\[
 A/\theta\longrightarrow\prod_{i\in I}A/\theta_i,\qquad
 [a]_\theta\longmapsto([a]_{\theta_i})_{i\in I}
\]
is a subdirect embedding.  The diagonal quotient map with domain
\(A\) is a subdirect embedding exactly when \(\theta=\Delta_A\).
These statements include \(I=\varnothing\), with
\(\theta=\nabla_A\) and \(A/\theta\) trivial.
\end{hbtheorem}

\begin{hbtheorem}{thm:p10-congruence-correspondence}{proved=annex:p10.subdirect.2}{Congruences above a quotient kernel}
For \(\theta\in\operatorname{Con}A\), the assignments
\[
\begin{aligned}
 \phi&\longmapsto
 \phi/\theta
   =\{([a]_\theta,[b]_\theta):(a,b)\in\phi\},
       &&\theta\subseteq\phi\in\operatorname{Con}A,\\
 \psi&\longmapsto
 \{(a,b):([a]_\theta,[b]_\theta)\in\psi\},
       &&\psi\in\operatorname{Con}(A/\theta)
\end{aligned}
\]
are inverse order isomorphisms between
\([\theta,\nabla_A]\) and \(\operatorname{Con}(A/\theta)\).
They preserve all meets and joins in these lattices and send
\(\theta\) to \(\Delta_{A/\theta}\).
\end{hbtheorem}

\begin{hbdefinition}{def:p10-subdirect-irreducible}{Subdirect irreducibility and the monolith}
An algebra \(A\) is \emph{subdirectly irreducible} when, for every
subdirect embedding
\[
 e:A\longrightarrow\prod_{i\in I}A_i
 \quad\text{with }I\ne\varnothing,
\]
at least one coordinate map \(\pi_i\circ e:A\to A_i\) is an
isomorphism.  The nonempty-index condition includes trivial algebras
in this convention; empty products remain allowed in
Definition~\ref{def:p10-subdirect-product}.

For a nontrivial algebra, a least congruence strictly above
\(\Delta_A\), when it exists, is called its \emph{monolith}.
\index{subdirectly irreducible algebra}\index{monolith}
\end{hbdefinition}

\begin{hbtheorem}{thm:p10-subdirect-monolith}{proved=annex:p10.subdirect.3}{Characterizing subdirect irreducibility}
An algebra \(A\) is subdirectly irreducible exactly when every
nonempty family \((\theta_i)_{i\in I}\) of congruences satisfying
\(\bigcap_i\theta_i=\Delta_A\) contains a member equal to
\(\Delta_A\).

Every trivial algebra has this property.  For nontrivial \(A\),
the following conditions are equivalent:
\[
\begin{gathered}
 A\text{ is subdirectly irreducible};\\
 \operatorname{Con}A\setminus\{\Delta_A\}
                 \text{ has a least member }\mu;\\
 \bigcap\bigl(\operatorname{Con}A\setminus\{\Delta_A\}\bigr)
                       \ne\Delta_A.
\end{gathered}
\]
In this case the intersection is the monolith \(\mu\).
For every \((a,b)\in\mu\) with \(a\ne b\),
\(\mu=\operatorname{Cg}_A(a,b)\).
\end{hbtheorem}

\begin{hbtheorem}{thm:p10-subdirect-representation}{proved=annex:p10.subdirect.4}{Birkhoff's subdirect representation theorem}
Assuming the axiom of choice, every nonempty finitary algebra \(A\)
is isomorphic to a subdirect product of subdirectly irreducible
homomorphic images of \(A\).

For nontrivial \(A\), one can take one quotient \(A/\theta_{a,b}\)
for every ordered pair \(a\ne b\), where \(\theta_{a,b}\) is maximal
among the congruences omitting \((a,b)\).  Then
\[
 a\longmapsto([a]_{\theta_{u,v}})_{u\ne v}
\]
is the required embedding.  A trivial algebra has a one-factor
representation by itself; alternatively it is isomorphic to the
empty product.

If \(A\) is finite and nontrivial, at most
\(|A|(|A|-1)\) finite quotient factors suffice.  This finite case
does not require the axiom of choice.
\end{hbtheorem}

The congruence-intersection criterion and the monolith characterization
correspond to Lemma~8.2 and Theorem~8.4 of the cited edition;
the representation theorem and its finite case are Theorem~8.6 and
Corollary~8.7 \cite[Chapter~II, \S8, pp.~57--59]{BurrisSankappanavar2012}.
The annex proves the chain-union and quotient steps explicitly.
Maximality here means maximal while separating a specified pair:
it need not mean maximal among all proper congruences.

\paragraph{A proper subdirect product.}
Let \(C_3=\{0<c<1\}\) be the three-element bounded lattice, and
\(C_2=\{0<1\}\) the two-element bounded lattice.
The maps with value lists
\[
 h_1=(0,0,1),\qquad h_2=(0,1,1)
\]
on \(0,c,1\) are monotone and preserve the endpoints.
Monotone maps between chains preserve binary minima and maxima,
so these maps preserve meets, joins, and the two constants.
Both are onto, and their kernels intersect in \(\Delta_{C_3}\).
They embed \(C_3\) as
\(\{(0,0),(0,1),(1,1)\}\leq C_2\times C_2\).
This is a three-element subdirect product inside a four-element
direct product.  Each factor is subdirectly irreducible because
its only congruences are equality and the universal relation.
